% Study note (not thesis text): the easy parts behind Section 4.1 (Generalised Linear Mixed Model),
% Thesis_coronary/content/technical/01_mixed_model.tex.
% Written 2026-11-16 on request ("for the easy stuff place it into a .tex for my understanding"), after
% the review that cut the elementary explanations from the thesis section.
% Sources: Dobson and Barnett, An Introduction to Generalized Linear Models (PDF in the repository
%   root; Sections 3.4, 5.7, 7.2, 7.3, 9.2, 10.2, 11.1, 11.5, 11.7) and Austin et al. 2017 (PMC5299617,
%   Section 2.2). Everything in a "My reasoning" box is algebra or interpretation that these two
%   sources do not state; do not cite it to them.
% All numbers are made up for illustration and were computed by hand on 2026-11-16.
\documentclass[11pt]{article}
\usepackage[a4paper,margin=25mm]{geometry}
\usepackage{amsmath,amssymb}
\usepackage{booktabs}
\usepackage[hidelinks]{hyperref}
\usepackage{parskip}
\usepackage{xcolor}
\usepackage{tcolorbox}
\newtcolorbox{thesisline}{colback=blue!4,colframe=blue!40!black,boxrule=0.4pt,left=4pt,right=4pt,top=3pt,bottom=3pt,title={Thesis sentence this explains},fonttitle=\small\bfseries}
\newtcolorbox{mine}{colback=orange!6,colframe=orange!60!black,boxrule=0.4pt,left=4pt,right=4pt,top=3pt,bottom=3pt,title={My reasoning (not in Dobson or Austin)},fonttitle=\small\bfseries}

\title{The mixed model of Section 4.1, the easy way\\[2pt]\large What each sentence means, with numbers}
\author{Study note, 16 November 2026}
\date{}

\begin{document}
\maketitle

This note follows the thesis section in order. Each part starts with the thesis sentence and then
explains it slowly. All numbers are invented.

In the thesis, a \emph{unit} is a coronary segment and a \emph{cluster} is a person. The technical
background says unit and cluster on purpose, since it describes the general method.

\section{Why a link is needed, and $g$ versus $g^{-1}$}

\begin{thesisline}
The two are therefore connected by a monotone link function $g$, whose inverse maps the linear
predictor to a risk: $p_i = g^{-1}(\eta_i)$.
\end{thesisline}

The linear predictor $\eta$ is a sum, so it can be any number. The risk $p$ must lie between 0 and 1.
The link is the translator between the two scales.

\begin{itemize}
  \item $g$ goes \textbf{up}: from the risk to the unbounded scale, $\eta = g(p)$.
  \item $g^{-1}$ goes \textbf{down}: from the linear predictor back to a risk, $p = g^{-1}(\eta)$.
\end{itemize}

The old sentence said that $g$ transforms the predictor into a number between 0 and 1. That is the
job of $g^{-1}$, and the equation printed $g^{-1}$, so text and equation disagreed. Monotone means
that a higher $\eta$ always gives a higher $p$, so the order of the units is kept.

\section{The hazard in three steps}

\begin{thesisline}
The link arises when the event is described through its hazard, which is the rate at which the event
occurs among the units that are still event-free. Summed over the window, the hazard gives the
cumulative hazard $\Lambda_i$, and the probability of no event by the end of the window is
$\exp(-\Lambda_i)$.
\end{thesisline}

\textbf{Step 1, the hazard.} Take segments without plaque. If 1\% of those still plaque-free develop
plaque each year, the hazard is 0.01 per year. It is a rate among those who have not yet had the
event.

\textbf{Step 2, the cumulative hazard.} Add the hazard up over the window. With 0.01 per year for
10 years, $\Lambda = 0.10$.

\textbf{Step 3, the probability of no event.} Staying event-free for 10 years means staying
event-free every single year, so the yearly probabilities multiply. A product of many small steps
turns into an exponential of the sum:
\[
  P(\text{no event}) = \exp(-\Lambda) = \exp(-0.10) = 0.905,
  \qquad
  p = 1 - 0.905 = 0.095 .
\]
The risk (0.095) is a little below the cumulative hazard (0.10). When events are rare the two are
close, which is why a cumulative hazard can loosely be read as a risk.

\section{Proportional hazards, and why the link does not ``assume'' it}

\begin{thesisline}
Under proportional hazards, a covariate multiplies the hazard by the same factor throughout the
window. The linear predictor is then the logarithm of the cumulative hazard, in which the intercept
holds the cumulative hazard of a unit whose covariates are all zero.
\end{thesisline}

Proportional hazards says: if a covariate raises the hazard by 50\%, it does so in year 1, in year 5
and in year 10 alike. Then the whole cumulative hazard is also 50\% higher:
\[
  \Lambda_i = \Lambda_0 \times \exp(\beta_1 x_{i1} + \beta_2 x_{i2} + \dots) .
\]
$\Lambda_0$ is the cumulative hazard of a unit with all covariates at zero, the baseline. Take the
logarithm:
\[
  \log \Lambda_i = \underbrace{\log \Lambda_0}_{\beta_0} + \beta_1 x_{i1} + \beta_2 x_{i2} + \dots = \eta_i .
\]
So the intercept $\beta_0$ \emph{is} the log baseline cumulative hazard. This was the unstated step
in the old text. Insert $\Lambda_i = \exp(\eta_i)$ into Step 3 and the cloglog risk appears:
\[
  p_i = 1 - \exp\{-\exp(\eta_i)\} .
\]

\begin{mine}
Why the old sentence ``the cloglog link assumes proportional hazards'' was not accurate: a link is
only a function. One can fit a cloglog model to any yes/no outcome and it will run. Proportional
hazards is the condition under which the coefficients may be \emph{read} as log hazard ratios.
Dobson introduces the cloglog link in Section 7.3 from a different story altogether (a dose-response
model), with no hazards in it. The route from grouped proportional hazards to the cloglog risk is
Prentice and Gloeckler 1978, which still needs to be opened (CHECK in the header).

With one follow-up scan, proportional hazards cannot be tested. That belongs in the Discussion.
\end{mine}

\section{The hazard ratio and its unit}

\begin{thesisline}
An increase of one in $x_{ik}$ with the other covariates held fixed multiplies the hazard by
$\exp(\beta_k)$, referred to as the hazard ratio. Since the other covariates are held fixed, the
hazard ratio is adjusted for them. A hazard ratio refers to one unit of its covariate and is read
together with that unit.
\end{thesisline}

If $\beta_k = 0.405$, then $\exp(0.405) = 1.5$: one step up in the covariate gives 1.5 times the
hazard. What ``one step'' is depends on how the covariate is scaled:

\begin{center}
\begin{tabular}{lll}
\toprule
Covariate scaled as & One step is & Hazard ratio 1.5 means \\
\midrule
per year of age & 1 year & 50\% higher hazard per year older \\
per mm & 1 mm & 50\% higher hazard per mm \\
per SD & 1 standard deviation & 50\% higher hazard per SD \\
\bottomrule
\end{tabular}
\end{center}

The same association gives very different numbers on different scales, so a hazard ratio without
its unit says nothing. The thesis reports the geometry hazard ratios per SD.

\textbf{Hazard ratio is not risk ratio.} With the numbers of Section 2: the hazard ratio 1.5 moves
$\Lambda$ from 0.10 to 0.15, and the risk from 0.095 to $1 - \exp(-0.15) = 0.139$. The risk ratio is
$0.139 / 0.095 = 1.46$, a little below 1.5.

\textbf{Adjusted.} ``Held fixed'' and ``adjusted'' are the same thing said twice: the comparison is
between two units that are alike in every other covariate of the model. The old text said it in two
subsections; now it is said once.

\begin{mine}
Why the sentence ``an adjusted hazard ratio different from 1 indicates an association'' was cut: an
\emph{estimated} hazard ratio is never exactly 1. 1.02 is ``different from 1'' and means nothing.
Whether there is an association is judged by the test and the confidence interval, not by the
estimate being unequal to 1. Adjusted for the six risk factors also does not mean adjusted for
everything, so it is an association and not a cause (root \texttt{CLAUDE.md}, prognostic factor
framing).
\end{mine}

\section{Exposure: offset or estimated $\gamma$}

\begin{thesisline}
The exposure enters the linear predictor as $\gamma \log E_i$, so it multiplies the cumulative
hazard by $E_i^{\gamma}$. When $\gamma$ is fixed at 1, the cumulative hazard is proportional to the
exposure and the term is referred to as an offset.
\end{thesisline}

Exposure is how much opportunity a unit has for the event: a longer follow-up, or a longer segment.
Since $\exp(\gamma \log E) = E^{\gamma}$, the term multiplies the cumulative hazard by $E^\gamma$.

\begin{center}
\begin{tabular}{lll}
\toprule
$\gamma$ & Doubling $E$ multiplies $\Lambda$ by & Reading \\
\midrule
1 (offset) & $2^1 = 2$ & twice the length, twice the hazard \\
0.5 (estimated) & $2^{0.5} = 1.41$ & hazard grows slower than the exposure \\
0 & $2^0 = 1$ & exposure does not matter \\
\bottomrule
\end{tabular}
\end{center}

An offset is a belief built into the model (hazard exactly proportional to exposure). Estimating
$\gamma$ lets the data decide, at the price of one parameter. Dobson defines the offset for Poisson
regression (Section 9.2), where the same $\log$ term appears; the estimated $\gamma$ is not in
Dobson.

\section{Why ``mixed'': fixed and random effects}

\begin{thesisline}
The coefficients $\beta_k$ are the same for all clusters and are referred to as fixed effects, while
$b_i$ differs between clusters, and a model with both is a mixed model.
\end{thesisline}

Nine segments of one person share that person's genes, diet and everything else that is not
measured. Their outcomes are therefore alike, and a model that treats 9 segments from 1,000 people
as 9,000 independent observations believes it has more information than it has.

\begin{itemize}
  \item \textbf{Fixed effects} $\beta_k$: one number each, the same for everybody (the effect of age,
        of curvature, \dots).
  \item \textbf{Random effect} $b_i$: one number per person, shifting all nine of that person's
        segments up or down together. The model does not report 1,000 separate $b_i$; it summarises
        them by their spread $\sigma^2$.
  \item \textbf{Mixed} = both kinds in one model. These are Dobson's words (Section 11.5).
\end{itemize}

\textbf{The index change.} Before Eq.~4.4, $i$ is a unit. From Eq.~4.4 on, $i$ is the cluster
(person) and $j$ the unit (segment) inside it, so $p_{ij}$ is the risk of segment $j$ in person $i$.
The thesis now says this in one sentence.

\textbf{What was weakened.} The old text said that ignoring clustering makes standard errors too
small between clusters and too large within. Dobson only shows that in one table as an example and
states no general rule, so the thesis now says what Section 11.1 does state: inference from such a
model is not valid.

\section{Frailty and the median hazard ratio}

\begin{thesisline}
Adding $b_i$ to the linear predictor multiplies the hazard of every unit in cluster $i$ by
$\exp(b_i)$, referred to as the frailty of the cluster. The MHR is the median ratio between the
hazards of two units with the same covariates from two randomly chosen clusters, with the higher
hazard divided by the lower.
\end{thesisline}

\textbf{Frailty.} $b_i = 0.4$ gives a frailty of $\exp(0.4) = 1.49$: every segment of this person
has about 1.5 times the hazard its covariates predict. $b_i = -0.4$ gives $0.67$. A frailty of 1 is
the typical person.

\textbf{The MHR, slowly.} $\sigma^2$ is hard to read, so it is translated to the hazard-ratio scale.
Pick two people at random. Take one segment from each, with identical covariates. Divide the higher
hazard by the lower. Do this for all possible pairs and take the median:
\[
  \mathrm{MHR} = \exp\!\left\{\sqrt{2\sigma^2}\;\Phi^{-1}(0.75)\right\} \approx \exp(0.954\,\sigma) .
\]

\begin{center}
\begin{tabular}{ll}
\toprule
$\sigma$ & MHR \\
\midrule
0 & 1.00 (people do not differ beyond their covariates) \\
0.5 & 1.61 \\
0.73 & 2.00 \\
1.0 & 2.60 \\
\bottomrule
\end{tabular}
\end{center}

``Median'' is why the thesis no longer says ``typically twice the hazard''. An MHR of 2 says: in
half of the pairs the two hazards differ by a factor of 2 or more, in the other half by less. The
MHR sits on the same scale as the covariate hazard ratios, so the two can be compared: an MHR of 2
next to a curvature hazard ratio of 1.2 per SD says that unmeasured differences between people
matter more than one SD of curvature. Austin's formula is for a Cox model with log-normal frailty;
the thesis model is a cloglog GLMM with the same normal $b_i$.

\section{Conditional hazard ratios}

\begin{thesisline}
The hazard ratio $\exp(\beta_k)$ is conditional on $b_i$, so it compares units with the same
frailty. For a covariate that varies between the units of a cluster, this is a comparison within a
cluster. For a covariate that is the same for all units of a cluster, it is a comparison between two
clusters with the same frailty.
\end{thesisline}

\begin{mine}
In the thesis this splits the covariates in two groups.
\begin{itemize}
  \item \textbf{Geometry} (curvature, taper per segment) varies within a person. Its hazard ratio
        can be read as: two segments of the same person, one SD apart in curvature.
  \item \textbf{The six risk factors} (and the LM angle and diameter ratios, if they are one value
        per person) are the same for all nine segments. Nobody has two ages. Their hazard ratio
        compares two people who are equally frail, which is a comparison nobody can observe
        directly.
\end{itemize}
A hazard ratio averaged over all people (a marginal one) is in general closer to 1 than the
conditional one, and more so the larger $\sigma^2$ is. The thesis has no opened source for this, so
the claim is not in the text. It is a point for the Discussion, with a source.

Root \texttt{CLAUDE.md} also notes that $\beta$ mixes between-person and within-person information
for the segment-level features. The clean ``within a person'' reading above is the reading of the
model, not a guarantee about where the information came from.
\end{mine}

\section{The likelihood ratio test, step by step}

\begin{thesisline}
The parameters of the GLMM are estimated by maximum likelihood. The random intercepts are not
observed, so the likelihood is averaged over their distribution. The model is fitted with and
without the set on the same units, and the two fits are compared through
$\mathrm{LR} = 2(\ell_1 - \ell_0)$.
\end{thesisline}

\textbf{Likelihood.} For given parameter values, how probable were the outcomes that were actually
seen? Maximum likelihood picks the parameters that make them most probable. $\ell$ is the logarithm
of that maximum. Since $b_i$ is unknown for every person, the program tries a handful of values of
$b_i$ and takes a weighted average (quadrature, in the appendix).

\textbf{The test.} Fit the model twice on exactly the same segments:
\begin{itemize}
  \item without geometry: $\ell_0 = -1250.0$
  \item with the 4 geometry covariates: $\ell_1 = -1243.0$
\end{itemize}
More covariates never fit worse, so $\ell_1 \geq \ell_0$ always. The question is whether the gain is
larger than chance would give.
\[
  \mathrm{LR} = 2 \times (-1243.0 - (-1250.0)) = 14.0, \qquad q = 4 .
\]
If geometry added nothing, LR would behave like a $\chi^2$ variable with 4 degrees of freedom, which
exceeds 9.49 only 5\% of the time. Here $14.0 > 9.49$, and $p = 0.007$. One test, for the whole set
at once: this is the main result of the thesis.

\textbf{Careful with $q$.} In Dobson Section 5.7, $q$ is the number of parameters of the smaller
model and the degrees of freedom are $p - q$. In the thesis $q$ is the number of added covariates.

\textbf{The caveat sentence.} ``This reference distribution assumes that the random effects capture
the dependence within clusters'': if segments of one artery are more alike than segments of
different arteries in the same person, the single person intercept does not capture all dependence,
and the $\chi^2$ reference is too optimistic.

\section{From segment hazards to a person risk, and dilution}

\begin{thesisline}
The cumulative hazard of a single unit is only one of the $n_i$ terms of $H_i$, so the association
of a covariate with the event in that unit is diluted in the cluster hazard.
\end{thesisline}

A person is event-free only if all nine segments are. Probabilities of ``no event'' multiply, and
with the cloglog form the product of exponentials is the exponential of a sum:
\[
  \exp(-\Lambda_{i1}) \times \dots \times \exp(-\Lambda_{i9}) = \exp\{-(\Lambda_{i1} + \dots + \Lambda_{i9})\} = \exp(-H_i) .
\]
(The frailty $e^{b_i}$ multiplies every term, so it comes out in front: $\exp(-e^{b_i} H_i)$.)

\textbf{Dilution with numbers.} Nine segments, each with $\Lambda = 0.02$, and frailty 1:

\begin{center}
\begin{tabular}{lll}
\toprule
 & Before & One segment gets hazard ratio 1.5 \\
\midrule
That segment's $\Lambda$ & 0.020 & 0.030 \\
That segment's risk & 1.98\% & 2.96\% \;(+49\%) \\
Person hazard $H_i$ & 0.180 & 0.190 \\
Person risk $1 - \exp(-H_i)$ & 16.5\% & 17.3\% \;(+5\%) \\
\bottomrule
\end{tabular}
\end{center}

A clear change in one segment barely moves the person risk. This is why the thesis judges geometry
by the LR test and the hazard ratios, and treats the change in C as secondary: a real association at
the segment level can leave the person-level ranking almost unchanged.

\textbf{``Before the frailty is applied.''} $\Lambda_{ij} = \exp(\eta_{ij})$ is the cumulative
hazard a segment would have in a person with frailty 1. The actual one is $e^{b_i}\Lambda_{ij}$.

\textbf{New person.} For somebody outside the data, $b_i$ is unknown, so the risk is averaged over
all possible frailties (appendix). Plugging in $b_i = 0$ instead gives the risk of the median
person, which is not the average risk.

\section{Questions an examiner could ask about this section}

\textbf{Why cloglog and not logistic regression?} Plaque is seen only at the follow-up scan, not
when it formed. Under proportional hazards the cloglog coefficients are log hazard ratios, and the
segment hazards add to a person hazard. The logit link has neither property. At low event rates
the two links fit almost equally well (appendix on the link choice).

\textbf{Can you check proportional hazards?} No, not with one follow-up scan. It is an assumption
of the reading, stated as such.

\textbf{What does your hazard ratio for curvature compare?} Two segments with the same risk factors
and the same person frailty, one SD apart in curvature. It is conditional on the random intercept.

\textbf{Why a random intercept and not robust standard errors?} The intercept gives a joint model
of the nine segments, which is needed to pool them into one person risk. Robust standard errors fix
the uncertainty only.

\textbf{Is one random intercept enough?} It assumes all nine segments of a person are equally
alike. Segments of the same artery may be more alike. A nested effect is described in the appendix
and is hard to estimate with few events.

\textbf{How well is $\sigma^2$ estimated?} Poorly when clusters are small and events rare, since
most information comes from people with more than one segment with plaque (appendix).

\textbf{Is the LR test valid in a mixed model?} Yes for fixed effects fitted by maximum likelihood
on the same data, given that the random-effects structure is right (Dobson 11.7). A test of a
variance being zero is different: the p-value is halved (appendix).

\textbf{Does a significant LR test mean geometry improves prediction?} No. See the dilution table.
The test says geometry carries information; how much it adds to the person risk is estimated
separately and may be small.

\end{document}
